% Appendix A body — builder emits \section{...}\label{sec:app:sky}; body only.
% Primary source: docs/derivations/G2a_completion_sky_marginal_4pi.md
This appendix derives the isotropic sky marginal that enters the completion
term of the likelihood (Sections~\ref{sec:framework} and \ref{sec:estimators}),
quantifies the flat-sky approximations behind it, and documents the one genuine
deviation uncovered by the derivation --- a missing $\sin\theta$ solid-angle
Jacobian --- which has since been corrected.

The pipeline approximates each event's GW likelihood by a trivariate Gaussian
in the coordinates $x = (\phi, \theta, u)$, where $(\phi,\theta)$ are the
ecliptic azimuth and polar angle and $u \equiv d_L/\hat d_L$ is the luminosity
distance in units of its point estimate $\hat d_L$:
\begin{equation}
  p(x_i \mid \phi, \theta, d_L)
  \simeq \mathcal{N}_3\big((\phi,\theta,u);\,\mu,\,\Sigma\big),
  \qquad \mu = (\hat\phi, \hat\theta, 1),
  \label{eq:app:sky:gauss}
\end{equation}
with $\Sigma$ assembled from the per-event Fisher (Cram\'er--Rao) covariance of
the 14-parameter EMRI model. Two conventions are load-bearing. First, the
Fisher matrix is computed with respect to the bare angles, so
equation~\eqref{eq:app:sky:gauss} is a probability density with respect to the
coordinate (Lebesgue) measure $\mathrm{d}\phi\,\mathrm{d}\theta\,\mathrm{d}u$,
\emph{not} the solid-angle measure
$\mathrm{d}\Omega = \sin\theta\,\mathrm{d}\theta\,\mathrm{d}\phi$. Second, the
in-catalogue numerator evaluates the same Gaussian pointwise at candidate-host
coordinates, so any sky \emph{integral} against an isotropic (per-steradian)
prior must carry the $\sin\theta$ Jacobian explicitly to remain consistent
with that branch.

With the uniform sky prior $p(\Omega) = 1/(4\pi)$ and the redshift prior
uniform in comoving volume--time \citep{Gray:2019ksv}, the completion-term
numerator is
\begin{equation}
  B_\mathrm{num}(h) = \int_{z_-}^{z_+} \mathrm{d}z\, \big[1 - f_k(z)\big]
  \left[ \int_{S^2} \frac{\mathrm{d}\Omega}{4\pi}\,
  \mathcal{N}_3\big((\phi,\theta,u(z,h));\,\mu,\Sigma\big) \right]
  \frac{w_\mathrm{pop}(z)}{4\pi},
  \label{eq:app:sky:bnum}
\end{equation}
where $f_k(z)$ is the catalogue completeness in the event's sky pixel,
$u(z,h) = d_L(z,h)/\hat d_L$, and
$w_\mathrm{pop}(z)/(4\pi) =
[\mathrm{d}V_\mathrm{c}/(\mathrm{d}z\,\mathrm{d}\Omega)]/(1+z)$ is the
per-steradian form of the full-sky population measure of
equation~\eqref{eq:wpop}, whose $1/(1+z)$ converts the source-frame
rate to the detector frame \citep{Mandel:2018mve}. Two distinct
factors of $1/(4\pi)$ thus appear: the isotropic sky-prior
normalization inside the square bracket, which is the factor at stake
in this appendix, and the constant measure conversion in
$w_\mathrm{pop}/(4\pi)$, which is common to $B_\mathrm{num}$, $D(h)$,
and $\beta_{\bar G}(h)$ and cancels in every ratio of
Section~\ref{sec:framework}.

Splitting $x = (\omega, u)$ with $\omega = (\phi,\theta)$, the Gaussian
factorizes exactly into a $u$-marginal times a sky conditional,
\begin{equation}
  \mathcal{N}_3\big((\omega,u);\,\mu,\Sigma\big)
  = \mathcal{N}_1\big(u;\,1,\,\Sigma_{uu}\big)\,
    \mathcal{N}_2\big(\omega;\,\mu_{\omega|u},\,\Sigma_{\omega|u}\big),
  \label{eq:app:sky:factor}
\end{equation}
where the first factor in equation~\eqref{eq:app:sky:factor} carries the
\emph{marginal} distance-fraction variance
$\Sigma_{uu} = \sigma_{d_L}^2/\hat d_L^{\,2}$ --- an element of the covariance,
not the smaller conditional variance $1/(\Sigma^{-1})_{uu}$. The distinction
matters: the pre-fix peak-density evaluation of the completion term
(Section~\ref{sec:pitfall}) effectively carried the too-narrow conditional
width on top of a $\sim 1/(2\pi\sigma_\phi\sigma_\theta)$ peak prefactor.
Since $\mathcal{N}_2$ integrates to unity under
$\mathrm{d}\phi\,\mathrm{d}\theta$, the sky integral in
equation~\eqref{eq:app:sky:bnum} reduces exactly to
$\bar p_\mathrm{GW}(u) = (4\pi)^{-1}\,\mathcal{N}_1(u;1,\Sigma_{uu})\,
\mathbb{E}_{\mathcal{N}_2}[\sin\theta]$. Extending the narrow sky Gaussian to
the plane gives
$\mathbb{E}_{\mathcal{N}_2}[\sin\theta] = \sin(\mu_{\theta|u})\,
\exp[-(\Sigma_{\omega|u})_{\theta\theta}/2]$, so the exact narrow-beam
isotropic sky marginal is
\begin{equation}
  \bar p_\mathrm{GW}(u)
  = \frac{\sin(\mu_{\theta|u})\,
    \mathrm{e}^{-(\Sigma_{\omega|u})_{\theta\theta}/2}}{4\pi}\,
    \mathcal{N}_1\big(u;\,1,\,\Sigma_{uu}\big)
  \approx \frac{\sin\hat\theta}{4\pi}\,
    \mathcal{N}_1\big(u;\,1,\,\Sigma_{uu}\big).
  \label{eq:app:sky:marginal}
\end{equation}
At the median sky-localization error of the simulated campaign,
$\Delta\Omega \approx 0.2\,\mathrm{deg}^2 = 6.1\times10^{-5}\,\mathrm{sr}$
($\sigma_\theta \approx \sigma_\phi \approx 3.3\times10^{-3}\,\mathrm{rad}$), % src: docs/derivations/G2a_completion_sky_marginal_4pi.md §4
the residual approximations are all negligible: truncating the Gaussian
support to the coordinate chart costs
$\mathcal{O}\big(\mathrm{e}^{-\hat\theta^2/2\sigma_\theta^2}\big)$; the
curvature of $\sin\theta$ across the beam enters through
$\tfrac12\sigma_\theta^2 \approx 5.6\times10^{-6}$; and the $u$-dependence of % src: docs/derivations/G2a_completion_sky_marginal_4pi.md §4
the conditional mean $\mu_{\theta|u}$ contributes a sub-percent, $h$-smooth
drift. The pull-out of the distance marginal is therefore valid to relative
error $\lesssim 10^{-5}$ ($\Delta\Omega/4\pi \sim 5\times10^{-6}$) --- % src: docs/derivations/G2a_completion_sky_marginal_4pi.md §4
\emph{provided} the $\sin\hat\theta$ factor is kept.

That factor is where the deviation occurred: the marginal was first
implemented with $\sin\hat\theta \to 1$. Because
equation~\eqref{eq:app:sky:gauss} is a coordinate density while the isotropic
prior is per-steradian, the exact marginal retains the Jacobian evaluated
under the beam; dropping it over-weights $B_\mathrm{num}$ relative to the
in-catalogue term $\beta_G L_\mathrm{cat}$ by $1/\sin\hat\theta$ per event ---
median $\approx 1.15$ and mean $\pi/2 \approx 1.57$ for isotropically % src: docs/derivations/G2a_completion_sky_marginal_4pi.md §4–§5
distributed events, unbounded only toward the ecliptic poles, where events are
in any case excluded by the Fisher condition-number gate. The factor is
$h$-independent event by event, but it multiplies only the completion term, so
it shifts the catalogue/completion mixture weight and thereby the shape of the
per-event likelihood; it does not cancel in the normalized posterior. Neither
dimensional analysis nor a re-reading of the Gaussian rescues it: interpreting
the sky Gaussian as a per-steradian density instead displaces the same
$\sin\theta$ inconsistency into the in-catalogue branch, and the
$\mathrm{rad}^2 \leftrightarrow \mathrm{sr}$ identification under which
$B_\mathrm{num}$ and $\beta_G L_\mathrm{cat}$ are dimensionally commensurable is
precisely where an $\mathcal{O}(1)$ dimensionless Jacobian hides. The
pixelated formulation of the completeness framework, by contrast, is immune by
construction: equal-area sky pixels implement the $\sin\theta$ measure exactly
\citep{Gray:2021sew}.
The omission is sub-leading --- three orders of magnitude below the
peak-density error of Section~\ref{sec:pitfall} --- and was found by this
derivation and corrected in the analysis code used for all results in this
paper. % src: .planning/gate/G7_systematics_budget.md row 4 (FIXED 4a259b7); docs/derivations/G2a_completion_sky_marginal_4pi.md §5

Equation~\eqref{eq:app:sky:marginal} passes four limiting cases. For a
complete catalogue, $f_k \to 1$, the integrand of
equation~\eqref{eq:app:sky:bnum} vanishes and the likelihood reduces to the
pure in-catalogue form of \citet{Gray:2019ksv}. For an empty catalogue,
$f_k \to 0$, $B_\mathrm{num}$ reduces to the catalogue-free volume-prior
likelihood --- the GW distance shell weighted by the comoving-volume prior. As
the sky localization sharpens, equation~\eqref{eq:app:sky:marginal} stays
finite and constant --- a perfectly localized event retains only the prior
weight $\sin\hat\theta\,\mathrm{d}\theta\,\mathrm{d}\phi/(4\pi)$ for an unseen
host under its beam --- whereas the peak-density evaluation it replaces
diverged as $1/(2\pi\sigma_\phi\sigma_\theta)$. Finally, for
$\sigma_{d_L} \to 0$, $\mathcal{N}_1 \to \delta(u-1)$ and the result remains
finite, the $\mathrm{d}z/\mathrm{d}u$ Jacobian emerging automatically because
the quadrature is carried out in $z$. The magnitude ratio between the replaced
and corrected evaluations, $4\pi/(2\pi\sigma_\phi\sigma_\theta) =
2/(\sigma_\phi\sigma_\theta)$, is $\approx 1.6\times10^{3}$ at a $2\degr$ sky
error and $\approx 1.8\times10^{5}$ at the median $0.2\,\mathrm{deg}^2$ % src: docs/derivations/G2a_completion_sky_marginal_4pi.md §7 item 5
localization: this is the $10^{3}$--$10^{5}$ over-weighting of the completion
term whose removal drives the de-railing of the posterior demonstrated in
Sections~\ref{sec:pitfall} and \ref{sec:realdata}.
